\documentclass[11pt]{article}
\usepackage[a4paper,margin=0.65in]{geometry}
\usepackage[T1]{fontenc}
\usepackage{amsmath,amssymb,amsfonts}
\usepackage{graphicx}
\usepackage{pdfpages}
\usepackage[english,shorthands=off]{babel}
\usepackage{fancyhdr}
\usepackage[bookmarks=false,colorlinks=false]{hyperref}
\usepackage[protrusion=true,expansion=false]{microtype}
\emergencystretch=3em
\tolerance=600
\providecommand{\modd}{\mathrm{mod}}
\providecommand{\Disc}{\mathrm{Disc}}
\providecommand{\canont}{}
\providecommand{\asterism}{*}
\providecommand{\Hessian}{H}
\providecommand{\tome}{}
\providecommand{\page}{p.}
\pagestyle{fancy}\fancyhf{}\fancyhead[L]{Pages 451--475}\fancyhead[R]{Cayley --- Collected Papers, Vol. XII}\renewcommand{\headrulewidth}{0.4pt}
\setlength{\headheight}{15pt}

\title{Collected Mathematical Papers, Vol. XII\\ \large Pages 431--455 (printed) / PDF pages 451--475\\ Papers 859 (concl.), 860, 861, 862, 863 (start)}
\author{Arthur Cayley}
\date{}

\begin{document}
\maketitle

\section*{859 (concluded). On the Lines Which Meet a Unicursal Quartic Curve.}

\noindent (Continuation of the determination of $\Omega$; printed page~431.)

\medskip

\noindent The expression thus becomes
\begin{align*}
\Omega &= P \cdot \quad x(y\gamma^3 - y^3\gamma) &&= \alpha\delta\Omega_1, \\
 &\quad + z(\alpha\delta^3 - \gamma^3\beta) &&= y(-z\beta_1 + 2w\theta\Omega) \\
 &\quad + 3xw(\alpha\gamma - \beta\delta) &&= 3\alpha\beta_1 \\
 &\quad + Q \cdot -3z(\beta\delta - w\gamma^2) &&= -2\alpha\delta\Omega_2 \\
 &\qquad + 3y(\alpha\delta - x\gamma^2) &&= 3\alpha\beta\delta_1 \\
 &\quad + R \cdot -3x(\beta^3 - z\alpha^2\gamma) &&= 3\alpha^2\delta_1 \\
 &\qquad - z(\alpha\delta - \beta^2) &&= (-\beta^4 + \alpha^3\delta) \\
 &\qquad + w(\alpha\delta^3 - \alpha\beta\gamma) &&= -\alpha^2\gamma\delta\Omega \\
 &\qquad + zw(-\alpha^2\delta^2 - \beta^2\gamma^2) &&= -2\alpha\beta\gamma\delta\Omega \\
 &\qquad + \tfrac{3}{2}\alpha\beta\delta - \tfrac{1}{2}\gamma^2 &&= \alpha\beta\delta\Omega \\
 &\quad + S \cdot xw(\alpha\beta\gamma - \beta\delta^2) &&= \alpha\beta\delta\Omega_2 \\
 &\qquad - z(\alpha\delta^2 - \alpha\beta\gamma) &&= -\alpha\beta\delta\Omega_1 \\
 &\qquad + yz(-2\alpha^2\beta\delta^2 - \beta^4\gamma^2) &&= -2\alpha\beta^2\delta\Omega \\
 &\qquad + \tfrac{3}{2}\alpha\beta\delta\gamma &&= 3\alpha\beta^2\delta\Omega \\
 &\qquad - \tfrac{1}{2}\beta^4 + 6\beta^2\gamma &&= -6\beta^3\Omega \\
 &\quad + w^2 \cdot 2\alpha^2\delta^3 - 4\alpha\beta\delta &&= -\alpha\beta^2\Omega \\
 &\qquad - (\alpha^3\beta\gamma - \beta^4) &&= -\beta^3\delta\Omega
\end{align*}

\medskip

\noindent and we thus finally obtain
\begin{align*}
\Omega \;=\; & PR\,(3\alpha\delta w - \gamma^2 y + S) \\
 & + RP\,(3\beta\gamma - \delta^2 z + \alpha w) \\
 & + PQ \cdot \tfrac{1}{3}\,S^2\, y \\
 & + Q\Omega \cdot - 3\,(\alpha w + \cdots).
\end{align*}

\medskip

\noindent viz.\ this is the equation of the cone, vertex $(\alpha,\beta,\gamma,\delta)$, which passes through the quartic curve $x:y:z:w = 1:\theta:\theta^2:\theta^3$. As regards the symmetry of this expression, it is to be remarked that, changing $(x,y,z,w)$ and $(\alpha,\beta,\gamma,\delta)$ into $(w,z,y,x)$ and $(\delta,\gamma,\beta,\alpha)$ respectively, we change $(\Theta,P,Q,R)$ and $(\Theta',P',Q',R')$ into $(\Theta,R,Q,P)$ and $(\Theta',R',Q',P')$ respectively, and so leave $\Omega$ unaltered. Again, interchanging $(x,y,z,w)$ and $(\alpha,\beta,\gamma,\delta)$, interchange $(\Theta,P,Q,R)$ and $(\Theta',P',Q',R')$, and so leave $\Omega$ unaltered.

\vfill
\newpage

\section*{860. On Briot and Bouquet's Theory of the Differential Equation $F\!\left(u,\,\dfrac{du}{dz}\right)=0$.}

\begin{flushright}\small
[From the \textit{Proceedings of the London Mathematical Society}, vol.~\textsc{xviii}.\ (1887), pp.~314--324.]
\end{flushright}

\medskip

\noindent\textsc{I consider} the theory of a differential equation of the first order as developed in Briot and Bouquet's \textit{Th\'eorie des fonctions doublement p\'eriodiques et en particulier des fonctions elliptiques} (\S 8, 1st ed., Paris, 1859), but I make some substantial variations in the mode of treatment.

I remark that, writing $u=x$, $\dfrac{du}{dz}=y$, I make the theory to depend altogether upon that of the curve $F(x,y)=0$; viz.\ the form of my result is, in order that a differential equation of the first order $F\!\left(u,\dfrac{du}{dz}\right)=0$ may have a monotropic integral, the curve $F(x,y)=0$ must satisfy certain conditions. Briot and Bouquet give their result (Theorem~iv., p.~296) under a somewhat different form, as follows; viz.\ changing the word ``monodrome'' into ``monotrope,'' their theorem is:

\medskip

\noindent\textit{``Pour qu'une \'equation diff\'erentielle du premier ordre de la forme
\[
\left(\frac{dy}{dz}\right)^{m} + f_{1}(u)\!\left(\frac{dy}{dz}\right)^{m-1} + \cdots + f_{m}(u) = 0
\]
admette une int\'egrale monotrope: $1^{\circ}$ les coefficients $f_{1}(u)$, $f_{2}(u)$, \ldots, $f_{m}(u)$ doivent \^etre des polyn\^omes entiers en $u$ et, au plus, le premier du second degr\'e, le second du quatri\`eme degr\'e, \ldots, le dernier du degr\'e $2m$; $2^{\circ}$ quand pour une certaine valeur de $u$ l'\'equation a une racine multiple diff\'erente de z\'ero, $\dfrac{dy}{dz}$ doit rester monotrope par rapport \`a $u$; $3^{\circ}$ quand pour une certaine valeur $u_{1}$ de $u$, l'\'equation a une racine multiple \'egale \`a z\'ero, le premier terme du d\'eveloppement de $\dfrac{dy}{dz}$ suivant les puissances croissantes de $(u-u_{1})^{1/n}$ doit avoir l'exposant $1-\dfrac{1}{n}$, si cet exposant est plus petit que l'unit\'e; $4^{\circ}$ enfin l'\'equation diff\'erentielle, que l'on d\'eduit de la premi\`ere en posant $u=\dfrac{1}{v}$, doit offrir pour $v=0$ les m\^emes caract\`eres. Ces conditions sont suffisantes.''}

\medskip

I notice that this may be regarded as a statement referring to the curve
\[
F(x,y) = y^{m} + f_{1}(x)\,y^{m-1} + \cdots + f_{m}(x) = 0,
\]
but that in $2^{\circ}$ a notion of monotropy is assumed, for $y$ must be a monotropic function of $x$; and that $4^{\circ}$ in effect introduces the new curve $F\!\left(\dfrac{1}{x},\,y\right)=0$.

\medskip

\noindent 1.\quad We have between the variables $u$ and $z$ a differential equation of the first order $F\!\left(u,\dfrac{du}{dz}\right)=0$, a rational and integral function of $u$ and $\dfrac{du}{dz}$. The equation determines $u$ as a function of $z$, and we wish to know whether $u$ is a monotropic function of $z$. It will not be so if we have a tropical point, viz.\ a point $z=c$, such that, in the neighbourhood thereof, the value of $u$ is given by a series $u-a = B(z-c)^{p/q} + \cdots$, in ascending powers of $z-c$, involving fractional powers of $z-c$. Conversely, if there is no tropical point, then $u$ will be a monotropic function of $z$; but we must consider and exclude not only tropical points at a finite distance, but also the point at infinity, viz.\ there must not be any development $u = Bz^{p/q} + \cdots$, in descending powers, involving fractional powers of $z$.

\medskip

\noindent 2.\quad A one-valued function (\textit{eindeutige Function}, or \textit{fonction uniforme} or \textit{bien d\'etermin\'ee}) is monotropic; since there is only one value, the function must, after a circuit described by the point $z$, recover its original value. But, conversely, a monotropic function is one-valued, that is, no two- or more-valued function can be monotropic. For, suppose a two-valued function of $z$, having the values $Z_{1}$ and $Z_{2}$, is monotropic; there is by hypothesis no tropical point, that is, no point $z=c$, such that the function $Z_{1}$, after a circuit described by the point $z$, instead of recovering its original value, assumes the value $Z_{2}$; and the like as to $Z_{2}$. But, this being so, it is meaningless to regard $Z_{1}$ and $Z_{2}$ as two values of the same function; they must be considered as two distinct and separate functions, each of them a one-valued monotropic function. And it thus appears that the two notions, monotropic and one-valued, are in fact equivalent.

\medskip

\noindent 3.\quad In what precedes, $u$ and $z$ are regarded as complex variables geometrically represented by means of two real points in their own planes respectively. But we now write $\dfrac{du}{dz} = y$, $u=x$, whence $F(x,y)=0$, regarding $x$, $y$ as the coordinates, real or imaginary, of a point of the curve represented by this equation. And this curve has, moreover, to be considered from a non-projective point of view; we have to distinguish between, and separately consider, finite points ($x$ and $y$ each of them finite), and the points at infinity ($x$ and $y$ either or each of them infinite). And it is moreover necessary to distinguish between, and separately consider, points on the axis of $x$ (or say axial points), and points off the axis of $x$ (say non-axial points). Taking, for convenience, $a$ to denote a finite value which may be $=0$, $b$ a finite value which is not $=0$ (or say, $a$ not $=\infty$, $b$ not $=\infty$ or $0$), we have thus the six cases, $1^{\circ}\;(a,b)$; $2^{\circ}\;(a,0)$; $3^{\circ}\;(a,\infty)$; $4^{\circ}\;(\infty,b)$; $5^{\circ}\;(\infty,0)$; $6^{\circ}\;(\infty,\infty)$. I regard also the axes of $x$ and $y$ as horizontal and vertical respectively, and speak of the tangent or element of the curve as horizontal, vertical, or inclined, accordingly.

\medskip

\noindent 4.\quad In the neighbourhood of any given point of the curve, we have $y=$ a series of powers of $x$, the form of the series being different for different classes of points. I write $P(x)$ to denote a series (finite or infinite) $A+Bx+Cx^{2}+\cdots$, in ascending powers of $x$; it is to be throughout understood that the leading coefficient $A$ is not $=0$. Of course, $P(x-a)$, $P(x-a)^{1/n}$, $P(x^{-1})$, \&c., will denote the like series $A+B(x-a)+\cdots$, $A+B(x-a)^{1/n}+\cdots$, $A+Bx^{-1}+\cdots$, \&c. The coefficients after the leading coefficient $A$ may any or all of them vanish; thus $P(x-a)^{1/n}$ extends to denote the series $P(x-a)$ of integer powers, but naturally the former notation is not used unless the series contains fractional powers. We can, by means of the symbol $P$, express for any given point of the curve the form of the expansion in the neighbourhood thereof; and, attaching to the point the expansion which belongs to it, we may, for instance, speak of a point $y=P(x-a)$, viz.\ if $b$ is the constant term of the series, then this is a non-axial point $(a,b)$, which is an ordinary (non-singular) point of the curve, and for which the element is not vertical; a like point with the element vertical would be a point $y=b+(x-a)^{\frac{1}{2}}P(x-a)^{\frac{1}{2}}$. Observe that, in the case of a multiple point, there is a separate expansion for each branch through the point, and in thus attaching an expansion to the point we regard the point as belonging to one of these branches; in dealing with a multiple point, it is necessary to consider separately all the different expansions, that is, all the branches through the point. It is clear that for a point $(a,b)$, $(a,0)$, or $(a,\infty)$ the expression for $y$ depends on $P(x-a)$, or $P(x-a)^{1/n}$; while, for a point $(\infty,b)$, $(\infty,0)$, or $(\infty,\infty)$, it depends on $P(x^{-1})$ or $P(x^{-1/n})$; viz.\ in the two cases respectively, we have a series in ascending powers of $x-a$, or in descending powers of $x$.

\medskip

\noindent 5.\quad In regard to any given point of the curve, substituting for $y$, $x$ their values $\dfrac{du}{dz}$, $u$, and thence forming the reciprocal of $\dfrac{du}{dz}$, we obtain $\dfrac{dz}{du} =$ a series in $u$; viz.\ this is a series in ascending powers of $(u-a)$, or in descending powers of $u$. Integrating, we have $z-c =$ a series in $u$, which series may or may not contain a logarithmic term $\log u$ or $\log(u-a)$; and, reverting the series, we obtain $u=$ a series in $(z-c)$. This result, applicable to the neighbourhood of the point $u=a$, may contain only integer powers of $z-c$, and be accordingly of the form $u=$ a one-valued series in $z$; and we then say that the point on the curve is a ``permissive'' point. Or it may contain fractional powers of $z-c$, and be accordingly of the form $u=$ a more-than-one-valued series in $z$; and we then say that the point on the curve is a ``prohibitive'' point. The necessary and sufficient condition in order that $u$ may be a monotropic function of $z$ is: the points of the curve must be all of them permissive points; or, what is the same thing, there must not be any prohibitive point.

\medskip

\noindent 6.\quad I form a complete table of permissive points, as follows:
\begin{center}
\begin{tabular}{ll}
$(a,b)$, & $y = P(x-a)$, \\
$(a,0)$, & $y = (x-a)P(x-a)$ or $(x-a)^{\frac{1}{n}}P(x-a)^{\frac{1}{n}}$, \\
$(a,\infty)$, & $y = (x-a)^{-\frac{1}{2}}P(x-a)^{\frac{1}{2}}$, \\
$(\infty,b)$, & $y = P(x^{-1})$, \\
$(\infty,0)$, $(\infty,\infty)$, & $y = xP(x^{-1})$ or $x^{2}P(x^{-1})$, \\
 & $y = x^{1+\varepsilon/n}P(x^{-1/n})$
\end{tabular}
\end{center}
\noindent (where $n$ is a positive integer $=2$ at least, $\varepsilon=0$, $1$, or $-1$,) viz.\ it is only a point $(a,b)$, $(a,0)$, $(\infty,b)$ or $(\infty,\infty)$ which may be permissive; a point $(a,\infty)$ or $(\infty,0)$ is prohibitive.

\medskip

\noindent 7.\quad A point $y = P(x-a)$ is permissive. We have
\[
\frac{du}{dz} = P(u-a), \qquad \frac{dz}{du} = \frac{1}{P(u-a)} = P(u-a),
\]
\[
z-c = (u-a)\,P(u-a),
\]
\[
u-a = (z-c)\,P(z-c),
\]
the required result. The several steps will be readily understood; the reciprocal of a series $P(u-a)$ is a series of the like form $P(u-a)$; integrating this in regard to $u$ as a series in $u-a$ (no constant being added on the right-hand side), we obtain a series $(u-a)\,P(u-a)$; and finally, $z-c$ being equal to this expression, we obtain, by reversion, a like expression $u-a = (z-c)\,P(z-c)$.

We might, for convenience, have written $x$, $u$, $z$ in place of $x-a$, $u-a$, $z-c$, respectively. The proof would have run
\[
y = P(x),\quad \frac{du}{dz} = P(u),\quad \frac{dz}{du} = P(u),\quad z = u\,P(u),\quad u = z\,P(z),
\]
meaning $u-a = (z-c)\,P(z-c)$, as above; and, in what follows, this is done accordingly.

It is worth while to make an instance in which the integration can be performed: say we consider the point $(0,1)$ of the curve $y = \dfrac{1+x+x^{2}}{1+2x}$; we have
\[
\frac{du}{dz} = \frac{1+u+u^{2}}{1+2u}, \qquad \frac{dz}{du} = \frac{1+2u}{1+u+u^{2}},
\]
\[
z = \log(1+u+u^{2}), \quad \text{or}\quad u+u^{2} = e^{z}-1,
\]
giving a series $u = z\,P(z)$, or the point is a permissive point, as in the general case. We have here $u+u^{2} = (e^{z}-1)$, viz.\ $u$ is not a monotropic function of $z$; but this is by reason of a prohibitive point $(-\tfrac{1}{2},\infty)$.

\medskip

\noindent 8.\quad A point $y = (x-a)\,P(x-a)$ is permissive.

Writing as above $u$, $z$ in place of $u-a$, $z-c$, we have
\[
\frac{du}{dz} = u\,P(u), \qquad \frac{dz}{du} = \frac{1}{u\,P(u)} = \frac{1}{u}\,P(u),\quad\text{or say}\quad \frac{dz}{du} = \frac{1}{u} + P(u),
\]
whence
\[
m z = \log u + u\,P(u),
\]
that is,
\[
u\,e^{u\,P(u)} = e^{mz}, \quad \text{or say}\quad u\,P(u) = e^{mz},
\]
whence
\[
u = e^{mz}\,P(e^{mz}), \quad \text{a one-valued series}.
\]

I take a particular instance, the point $(0,0)$ for the curve $y = \dfrac{x - x^{2}}{1 + x^{2}}$; here
\[
\frac{du}{dz} = \frac{u - u^{2}}{1 + u^{2}}, \qquad \frac{dz}{du} = \frac{1+u^{2}}{u - u^{2}},
\]
\[
z = \log\!\frac{u}{1-u} - \tfrac{2}{3}u^{3}, \quad\text{or}\quad \frac{u}{1-u}\,e^{-\frac{2}{3}u^{3}} = e^{z};
\]
the point is permissive. The equation just obtained shows that $u$ is not a monotropic function of $z$; but this is by reason of the prohibitive points $(1,\infty)$.

\medskip

\noindent 9.\quad A point $y = (x-a)^{\frac{1}{2}}\,P(x-a)^{\frac{1}{2}}$ is permissive. Here
\[
\frac{du}{dz} = u^{\frac{1}{2}}\,P(u^{\frac{1}{2}}), \qquad \frac{dz}{du} = u^{-\frac{1}{2}}\,P(u^{\frac{1}{2}}),
\]
whence
\[
z = -\frac{A}{u^{\frac{1}{2}}} + B\log u + u\,P(u),
\]
where the logarithmic term occurs or does not occur, according as $B$ is not or is $=0$. In the former case, we have
\[
e^{z/B}\,u^{\frac{1}{2}} = e^{?};
\]
in the latter case,
\[
z = \frac{?}{u^{?}};
\]
and in each case we can, by reversion, express $u$ as a one-valued series in $z$. We may take as examples the two curves $y = \dfrac{1-x^{2}}{x^{\frac{1}{2}}}$ and $y = \dfrac{x^{\frac{1}{2}}}{1-x^{2}}$; in each case the point $(0,0)$ is a permissive point, but, by reason of the prohibitive points $(1,\infty)$ and $(-1,\infty)$, $u$ is not a monotropic function of $z$.

\medskip

\noindent 10.\quad As to the other forms of permissive points, the proofs depend upon those for the forms already considered. Thus for
\[
y = (x-a)^{\frac{1}{n}}\,P(x-a)^{\frac{1}{n}} \quad\text{or}\quad y = (x-a)^{1+\varepsilon/n}\,P(x-a)^{\frac{1}{n}},
\]
writing $v = u^{1/n}$, we find
\[
\frac{du}{dz} = v\,P(v) \cdot \text{(or)} \cdot v^{1+\varepsilon}\,P(v),
\]
that is,
\[
\frac{dv}{dz} = P(v), \quad v\,P(v),\quad \text{or}\quad v^{\frac{1}{2}}\,P(v);
\]
in each case $v=$ one-valued series in $z$, and thence $u=v^{n}=$ one-valued series in $z$.

Similarly for $y = P(x^{-1})$, $y = x\,P(x^{-1})$, $y = x^{2}\,P(x^{-1})$, or say here
\[
y = x^{1+\varepsilon}\,P(x^{-1});
\]
\[
\frac{du}{dz} = u^{1+\varepsilon}\,P(u^{-1}),
\]
or, writing $v = u^{-1}$, we have
\[
\frac{du}{dz} = u^{\varepsilon-1}\,P(u^{-1}) = v^{1-\varepsilon}\,P(v);
\]
viz.\ this is
\[
\frac{dv}{dz} = P(v), \quad v\,P(v), \quad\text{and}\quad v^{\frac{1}{2}}\,P(v);
\]
in each case $v=$ one-valued series in $z$, and thence $u = v^{-1} =$ one-valued series in $z$.

And finally, for
\[
y = x^{1+\varepsilon/n}\,P(x^{-1/n});
\]
here
\[
\frac{du}{dz} = u^{1+\varepsilon/n}\,P(u^{-1/n}),
\]
or, writing $v = u^{-1/n}$, we have
\[
\frac{dv}{dz} = v^{-(n-1)\varepsilon/n}\,P(v) \quad\text{or}\quad v^{?}\,P(v),
\]
that is,
\[
\frac{dv}{dz} = P(v),\quad v\,P(v),\quad\text{or}\quad v^{\frac{1}{2}}\,P(v);
\]
as before; $v=$ one-valued series in $z$, and thence $u = v^{-n} =$ one-valued series in $z$.

\medskip

\noindent 11.\quad There is no difficulty in showing that every point of the curve, not being a permissive point as above, is a prohibitive point. Thus for a point $(a,b)$, if the series for $y-b$ contain any fractional power, say, if we have $y = b + (x-a)^{1/n} + \cdots$, then
\[
\frac{du}{dz} = b + u^{1/n}, \qquad \frac{dz}{du} = \frac{1}{b} - \frac{u^{1/n}}{b^{2}} + \cdots,
\]
whence, reverting, $u = bz +$ series containing the fractional power $z^{1/n+1}$.

Again, a point $y = (x-a)^{m}\,P(x-a)$, $m=3$ or any higher integer value, is prohibitive. Attending only to the leading term, we have $\dfrac{du}{dz} = u^{m}$, whence $\dfrac{dz}{du} = u^{-m}$ or $u^{1-m}$, whence, reverting, $z = u^{1-m}$, $u = z^{1/(1-m)}$, which is a fractional power; and this also is the case if $m$ be a negative integer. And similarly, in other cases, the series for $u$ will always contain fractional powers of $z$.

\medskip

\noindent 12.\quad In order that $u$ may be a monotropic function of $z$, the points of the curve must be all of them permissive, or, what is the same thing, there must be no prohibitive point; we have to consider the geometrical signification of this condition. If we attend first to the ordinary (or non-singular) points of the curve, a non-axial point $(a,b)$ must be a point $y = P(x-a)$, viz.\ there must be no such point with a vertical element.

An axial point $(a,0)$ may be $y = (x-a)\,P(x-a)$, viz.\ the element may be inclined; or the point may be $y = (x-a)^{2}\,P(x-a)$, viz.\ the element may be horizontal (but in this case there must be an ordinary or two-pointic contact only); or the point may be $y = (x-a)^{\frac{1}{2}}\,P(x-a)^{\frac{1}{2}}$, viz.\ the element may be vertical.

The point must not be $(a,\infty)$, viz.\ there must be no asymptote parallel to (or coincident with) the axis of $y$.

The point may be $(\infty,b)$, $y = P(x^{-1})$, that is, there may be an ordinary, or osculating, asymptote parallel to the axis of $x$. But there is no point $(\infty,0)$, that is, there must be no asymptote coincident with the axis of $x$.

There may be points $(\infty,\infty)$, $y = x\,P(x^{-1})$, that is, ordinary or osculating asymptotes inclined to the axes; and there may also be points $(\infty,\infty)$, $y = x^{2}\,P(x^{-1})$, that is, asymptotic parabolas of vertical axis.

\medskip

\noindent 13.\quad We have, moreover, conditions as to the singular points; thus, every non-axial multiple point $(a,b)$ must be a point with each branch of the form $y = P(x-a)$, that is, each branch must be an ordinary (non-cuspidal) branch, and must be non-vertical. The axial multiple or singular points $(a,0)$ may have ordinary (non-cuspidal) branches $y = (x-a)\,P(x-a)$ and $y = (x-a)^{2}\,P(x-a)$, inclined or horizontal, or else $y = (x-a)^{\frac{1}{2}}\,P(x-a)^{\frac{1}{2}}$, vertical; and there may also be cuspidal branches $y = (x-a)\,P(x-a)^{1/n}$, $y = (x-a)^{1/n}$, and $y = (x-a)^{1+1/n}\,P(x-a)^{1/n}$, $P(x-a)^{1/n}$ inclined, horizontal, or vertical.

There must be no singular points $(a,\infty)$; any such point is, in fact, excluded by the condition, no asymptote parallel to or coincident with the axis of $y$.

There may be multiple points $(\infty,b)$, but not $(\infty,0)$; viz.\ as above, there may be asymptotes parallel to, but not coincident with, the axis of $x$; but in this case, the several branches must be each of them an ordinary branch $y = P(x^{-1})$.

Finally, there may be multiple or singular points $(\infty,\infty)$, but each branch must be either an ordinary branch $y = x\,P(x^{-1})$ or $y = x^{2}\,P(x^{-1})$, or a cuspidal branch.

\medskip

\noindent 14.\quad The enumeration seems long and difficult. But for a given curve we have, in fact, only to see that there are no ordinary non-axial points $(a,b)$ of vertical element, and to further examine the finite singular points $(a,b)$ and $(a,0)$, and also the several infinite branches as giving ordinary or else singular points $(a,\infty)$, $(\infty,b)$, $(\infty,0)$, or $(\infty,\infty)$; viz.\ it has to be seen for each infinite branch that the expansion is of the proper form.

\medskip

\noindent 15.\quad Writing the equation of the curve in the form
\[
f_{0}(x)\,y^{m} + f_{1}(x)\,y^{m-1} + \cdots + f_{m}(x) = 0,
\]
where $f_{0}$, $f_{1}$, \ldots, $f_{m}$ denote rational and integral functions of $x$, then, if $a$ be a root of the equation $f_{0}(x) = 0$, there will be on the curve a point $(a,\infty)$ which is prohibitive; $f_{0}(x)$ is therefore a constant, or, taking it to be $=1$, the form must be
\[
y^{m} + f_{1}(x)\,y^{m-1} + \cdots + f_{m}(x) = 0,
\]
where $f_{1}$, $f_{2}$, \ldots, $f_{m}$ denote rational and integral functions of $x$.

It is to be shown that the degrees of these functions are equal respectively to $2$, $4$, \ldots, $m$ at most. Supposing that the degrees are $2$, $4$, \ldots, $2m$; then for the points at infinity we find an expansion $y = Ax^{n} + \cdots$, in descending powers of $x$, by substitution in an equation
\[
y^{m} + a_{1}x\,y^{m-1} + a_{2}x^{2}\,y^{m-2} + \cdots + a_{m}x^{m} = 0;
\]
whence $n = 2$, and $A$ is determined by an equation of the order $m$; there are thus $m$ branches $y = x^{2}\,P(x^{-1})$, corresponding to permissive points. If, however, any one of the $y$ functions be of a higher degree, then it may be shown (but the formal proof is not easily given) that there will be at any rate a branch $y = Ax^{n} + \&c.$, where $n$ has a value $>2$, and we have thus a prohibitive point; hence the degrees must be as stated.

\medskip

\noindent 16.\quad As an example of a differential equation with a monotropic solution, take
\[
\left(\frac{du}{dz}\right)^{2} = \left(\frac{du}{dz}\right) + 27u^{2} - 4 = 0;
\]
we have here the curve
\[
y^{3} = 3y + 27x^{2} - 4 = 0,
\]
which is a nodal cubic, as shown in the annexed figure. To put the node in evidence, the equation may be written in the form $(y+2)^{2}(y-1) + 27x^{2} = 0$; viz.\ the node is the point $(0,-2)$, and the equation of the tangents is $(y+2)^{2} - 27x^{2} = 0$. The curve besides meets the line $x=0$ in the point $y=1$, where the tangent is horizontal, and it meets the axis $y=0$ in the points $x = \pm \dfrac{2}{3\sqrt{3}}$, at each of which the tangent is vertical; these two points, as lying on the axis $y=0$, are thus each of them permissive; and they are the only points where the tangent is vertical (in fact, the tangents from the point at infinity on the line $x=0$ are the two lines $x = \pm \dfrac{2}{3\sqrt{3}}$, and the line infinity counting twice, as a tangent at an inflexion).

\textit{[Figure: nodal cubic curve $y^{3} - 3y + 27x^{2} - 4 = 0$, showing the loop at the node $(0,-2)$ with branches descending and the upper smooth arc through $(0,1)$.]}

For the points at infinity, we have $y = -3x^{\frac{2}{3}} - 1 - \tfrac{1}{3}x^{-\frac{2}{3}} + \&c.$, which is of the form $y = x^{\frac{2}{3}}P\!\left(\dfrac{1}{x^{1/3}}\right)$, included in $y = x^{1+\varepsilon/n}\,P(x^{-1/n})$, and the point is thus permissive; there is no prohibitive point, and the differential equation has a monotropic solution. This is, in fact, the rational solution $u = (z-c) - (z-c)^{3}$, or say $u = z - z^{3}$; the curve is thus given by the two equations $x = z - z^{3}$, $y = 1 - 3z^{2}$.

\medskip

\noindent 17.\quad As a second example, take the differential equation
\[
\left(\frac{du}{dz}\right)^{3} - \left(\frac{du}{dz}\right)^{2} + 27u^{2} - 4 = 0.
\]
We have here the curve
\[
y^{3} - y^{2} + 4x^{2} - 27x^{4} = 0,
\]
which is a trinodal quartic curve, as shown in the figure. There is a node at the origin with the tangents $y^{2} - 4x^{2} = 0$, and, writing the equation in the form
\[
(3y+1)^{2}(3y-2)^{2} - 8(27x^{2}-2)^{2} = 0,
\]
we have for the other two nodes $x = \pm \dfrac{1}{3\sqrt{3}}$, $y = \dfrac{2}{3}$. The lowest points are given by $x = \pm \dfrac{2}{3\sqrt{3}}$, $y = -\dfrac{1}{3}$; viz.\ the line $y = -\dfrac{1}{3}$ is a horizontal tangent. Writing $x=0$, we have $y^{3} - y^{2} = 0$, that is, $y^{2} = 0$, the node at the origin, and $y=1$, the height of the loop; and, writing $y=0$, we have $4x^{2} - 27x^{4} = 0$, that is, $x^{2} = 0$, the node at the origin, and $x = \pm \dfrac{2}{3\sqrt{3}}$, the other two intersections with the axis of $x$. The tangents at these two points are vertical; as being on the axis, they are thus permissive points. And they are the only points with a vertical tangent; in fact, the point $(x=0, y=\infty)$ is a point of the curve, with the line $\infty$ as an osculating tangent (4-pointic intersection); hence the tangents from the point $x=0$, $y=\infty$ are the line $\infty$ counting four times, and the lines $x = \pm \dfrac{2}{3\sqrt{3}}$, touching at the points $\left(\pm \dfrac{2}{3\sqrt{3}},\, -\dfrac{1}{3}\right)$, as above.

\textit{[Figure: trinodal quartic curve $y^{3} - y^{2} + 4x^{2} - 27x^{4} = 0$, showing the three nodes and the lower loops with horizontal tangent at $y=-\tfrac{1}{3}$.]}

For the infinite branches, we have $y = 3^{\frac{2}{3}}x^{\frac{4}{3}} + \cdots$, which is of the right form $y = x^{1+\frac{1}{3}}P\!\left(\dfrac{1}{x^{1/3}}\right)$. We thus see that there is no prohibitive point; and the differential equation has a monotropic solution accordingly.

Writing $z$ in place of $z-c$, the solution in fact is
\[
u = \frac{e^{z} + e^{-z}}{2}\cdot\,?
\]
hence, putting $\theta = e^{z} - e^{-z}$, the curve is given by the two equations
\[
x = \theta - \theta^{3}, \qquad y = (1 - 3\theta^{2})\cdot;
\]
the monotropic function $u$ satisfies the differential equation.

\medskip

\noindent 18.\quad $F\!\left(u,\,\dfrac{du}{dz}\right)$ must be either a rational function, a singly periodic function, or a doubly periodic function of $z$; say the forms are
\[
u = P(z),\quad u = P(e^{z}),
\]
and
\[
u = P\{\operatorname{sn}(gz,k)\} + Q\{\operatorname{sn}(gz,k)\}\,\operatorname{cn}(gz,k)\,\operatorname{dn}(gz,k),
\]
where $P$, $Q$ denote rational functions. I do not at present consider the criteria (such as are given in Briot and Bouquet's Theorem~v., p.~301) for determining by means of the curve which is the form of the integral. I remark, however, that in the first and second cases the curve is unicursal, while in the third case it is bicursal; or say that, according as the deficiency is $0$ or $1$, the integral is rational or simply periodic, or else it is doubly periodic. Moreover, the curve being unicursal, we can express the coordinates as equal to rational functions $P(\theta)$, $Q(\theta)$ of a parameter $\theta$; and, being bicursal, we can express them as functions of the form
\[
P(\theta)\sqrt{1-\theta^{2}\cdot 1-k^{2}\theta^{2}} + P_{1}(\theta),\quad Q(\theta)\sqrt{1-\theta^{2}\cdot 1-k^{2}\theta^{2}} + Q_{1}(\theta).
\]
of the parameter $\theta$; and supposing the coordinates $(x,y)$, that is, $u$ and $\dfrac{du}{dz}$, thus expressed, it should be easy to establish the relation of $\theta$ with $z$, $e^{z}$ or $\operatorname{sn}(gz,k)$ in the three cases respectively.

\vfill
\newpage

\section*{861. Note on a Formula Relating to the Zero-Value of a Theta-Function.}

\begin{flushright}\small
[From \textit{Crelle's Journal der Mathem.}, t.~c.\ (1887), pp.~87, 88.]
\end{flushright}

\medskip

\noindent\textsc{I had} some difficulty in verifying for the case of a single theta-function, a formula given in Herr Thomae's paper ``Beitrag zur Theorie der $\vartheta$-Functionen,'' \textit{Crelle's Journal}, vol.~\textsc{lxxi}.\ (1870), pp.~201--222. The formula in question (see p.~216) is given as follows:
\[
(11)\qquad \vartheta(0,0,\ldots,0) = \sqrt{\frac{|A_{i,k}|}{(2\pi i)^{p}}}\,\sqrt[4]{\Disc(\theta,0,\ldots,0)\,\Disc(0,\theta,\ldots,0)\cdots},
\]
but the denominator factor should I think be $(\pi i)^{p}$ instead of $(2\pi i)^{p}$. Making this alteration, then in the case of a single theta-function, $p=1$, and the function belongs to the radical
\[
\sqrt{x-k_{1}\cdot x-k_{2}\cdot x-k_{3}\cdot x-k_{4}}
\]
where
\[
(k_{1},k_{2},k_{3},k_{4}) = \left(-\tfrac{1}{k},\,-1,\,+1,\,+\tfrac{1}{k}\right).
\]
The determinant $|A_{i,k}|$ is a single term $=A$, and the formula becomes
\[
\vartheta(0) = \sqrt{\frac{A}{\pi i}}\,\sqrt[4]{(k_{1}-k_{2})(k_{4}-k_{3})},
\]
where $k_{1}-k_{2}$, $k_{4}-k_{3}$ are each $= -1 + \dfrac{1}{k}$, and we have therefore
\[
\sqrt[4]{(k_{1}-k_{2})(k_{4}-k_{3})} = \sqrt{-1 + \frac{1}{k}} = \sqrt{\frac{1-k}{k}}.
\]
$A$ also denotes the integral
\[
\int_{k_{2}}^{k_{3}} \frac{dx}{\sqrt{x-k_{1}\cdot x-k_{2}\cdot x-k_{3}\cdot x-k_{4}}} \;=\; \int_{-1}^{+1} \frac{k\,dx}{\sqrt{1-x^{2}\cdot 1-k^{2}x^{2}}} = 2ikK,
\]
and the formula thus is
\[
\vartheta(0) = \sqrt{\frac{2ikK}{\pi i}\cdot\sqrt{\frac{1-k}{k}}} \;=\; \sqrt{\frac{2kK}{\pi}}\cdot \sqrt[4]{\frac{1-k}{k}}.
\]

But observe that, in the theta-function as defined by the equation
\[
\vartheta(x) = \Sigma e^{i\pi\omega n^{2}}\cdot,
\]
$\omega$ is used to denote the value
\[
\omega = a/B, \quad a = \tfrac{2\pi}{B}\,?
\]
where $A$ is the above-mentioned integral, and $B$ is the integral
\[
B = \int_{k_{1}}^{k_{2}} \frac{dx}{\sqrt{x-k_{1}\cdot x-k_{2}\cdot x-k_{3}\cdot x-k_{4}}} \;=\; \int_{-?}^{?} \frac{k\,dx}{\sqrt{1-x^{2}\cdot 1-k^{2}x^{2}}} = 2kK',
\]
which value must however be taken negatively, viz.\ we must write $B = -2kK'$, and we then have
\[
\omega = \frac{2ikK}{-2kK'} = -i\,\frac{K}{K'};
\]
viz.\ writing as usual
\[
e^{-\pi K'/K} = q,
\]
the $\omega$ of the theta-function is not $q$, but it is $r^{2}$; and the zero-value $\vartheta(0)$ is $= 1 + 2r^{2} + 2r^{8} + 2r^{18} + \cdots$. The equation thus is
\[
1 + 2r^{2} + 2r^{8} + \cdots \;=\; \sqrt{\frac{2kK}{\pi}}\cdot \sqrt[4]{\frac{1-k}{k}},
\]
which is right; in fact, writing $k'$ in place of $k$, and consequently $K'$, $q$ in place of $K$, $r$ respectively, the equation becomes
\[
1 + 2q^{2} + 2q^{8} + \cdots \;=\; \sqrt{\frac{2k'K'}{\pi}}\cdot \sqrt[4]{\frac{1-k'}{k'}};
\]
we have
\[
1 + 2q + 2q^{4} + \cdots \;=\; \sqrt{\frac{2K}{\pi}},
\]
and changing $q$ into $q^{2}$, then (\textit{Fund.\ Nova}, p.~92) $K$ is changed into $\tfrac{1}{2}K(1+k)$, and $k$ is changed into $\dfrac{1-k'}{1+k'}$; for small values of $q$, observe that we have
\[
\frac{2K}{\pi} = 1 + 4q + \&c.,
\]
and thence
\[
\frac{K(1+k)}{\pi} = 1 + 4q^{2} + \&c.,\quad\text{or}\quad k\sqrt{\frac{1+k}{2}} \;=\; \frac{1-k'}{1+k'}\,?
\]
as a verification of $q$, observe that we have
\[
\frac{2K}{\pi}\quad\text{and thence}\quad k = \frac{4\sqrt{q}}{1+\cdots},\quad\text{or}\quad k\sqrt{k} = ?
\]
we have the formula in question.

\medskip

\begin{flushright}\textit{Cambridge, 12 February, 1886.}\end{flushright}

\vfill
\newpage

\section*{862. Note on the Theory of Linear Differential Equations.}

\begin{flushright}\small
[From \textit{Crelle's Journal der Mathem.}, t.~c.\ (1887), pp.~286--295.]
\end{flushright}

\medskip

\noindent 1.\quad I consider a linear differential equation
\[
p_{0}\frac{d^{m}y}{dx^{m}} + p_{1}\frac{d^{m-1}y}{dx^{m-1}} + \cdots + p_{m}y = 0,
\]
where $p_{0}$, $p_{1}$, \ldots, $p_{m}$ are rational and integral functions of $x$, having no common factor: as usual, $x$ is a complex magnitude represented by a point, and we consider the integrals belonging to a singular point $x = a$ of the differential equation. An integral may be a regular integral, or it may be what Thom\'e calls a normal elementary integral: the theory of these integrals (which I would rather call subregular integrals) requires, I think, further examination.

\medskip

\noindent 2.\quad I retain $x$ as the independent variable, but for shortness use $t$ to denote $\dfrac{1}{x-a}$: and I take as the dependent variable $z$, $= \dfrac{1}{y}\,\dfrac{dy}{dx}$: we have then $z$ determined from any value of $z$ we derive a corresponding value $e^{\int z\,dx}$ of $y$.

\medskip

\noindent 3.\quad To obtain the $z$-equation, using for shortness accents to denote differentiation in regard to $x$, we have $z = \dfrac{y'}{y}$, and thence $z' = \dfrac{y''}{y} - \left(\dfrac{y'}{y}\right)^{2}$, that is, $\dfrac{y''}{y} = z' + z^{2}$; also
\[
\frac{y'''}{y} - \frac{y''\,y'}{y^{2}} = z'' + 2zz',
\]
that is
\[
\frac{y'''}{y} = z'' + 3zz' + z^{3},\quad\text{or finally}\quad z'' + 3zz' + z^{3} = \frac{y'''}{y},
\]
and so on; viz.\ the values of $\dfrac{y'}{y}$, $\dfrac{y''}{y}$, $\dfrac{y'''}{y}$, \ldots are $z$, $z^{2} + z'$, $z^{3} + 3zz' + z''$, \ldots; $\dfrac{y^{(n)}}{y}$, generally for the first term is $z^{n}$ and the last term is $z^{(n-1)}$. If, to fix the ideas, the $y$-equation is
\[
\frac{d^{3}y}{dx^{3}} + p_{1}\frac{d^{2}y}{dx^{2}} + p_{2}\frac{dy}{dx} + p_{3}y = 0,
\]
then, dividing by $y$, the $z$-equation is
\[
(z'' + 3zz' + z^{3}) + p_{1}(z' + z^{2}) + p_{2}\,z + p_{3} = 0;
\]
and similarly, from the $y$-equation of the order $m$, we derive a $z$-equation of the order $m-1$,
\[
p_{0}\,[z^{(m-1)} + \cdots + z^{m}] + p_{1}[\cdots + z^{m-1}] + \cdots + p_{m-2}(z' + z^{2}) + p_{m-1}\,z + p_{m} = 0.
\]

\medskip

\noindent 4.\quad In this equation, each coefficient is in the first instance a rational and integral function of $x$, or, what is the same thing, of $x-a$; writing $t$ to denote $\dfrac{1}{x-a}$, each coefficient is thus the sum of a finite number of negative integer powers of $t$; hence multiplying the whole equation by a proper positive power of $t$, the several coefficients become each of them the sum of a finite number of positive powers of $t$, or arranging in descending powers of $t$, say the general form is $p_{r} = \alpha_{r}t^{\lambda} + \beta_{r}t^{\lambda-1} + \cdots + \kappa_{r}$, where $\lambda$ is a positive integer which may be $=0$, and which has for each coefficient its proper value. We wish to satisfy the $z$-equation by a value $z = $ descending series in $t$, viz.\ the form is $z = At^{\alpha} + A't^{\alpha-1} + \cdots$, with powers which are integral or fractional, but where the number of positive powers is finite. The theory is almost identical with that of the solution in like form of the algebraical equation
\[
p_{0}\,z^{m} + p_{1}\,z^{m-1} + \cdots + p_{m-1}\,z + p_{m} = 0,
\]
or say the equation $U = 0$, where $U$ is a rational and integral function of $z$ and $t$, of the degree $m$ as regards $z$.

\medskip

\noindent 5.\quad Considering $z$ and $t$ as Cartesian coordinates, the equation in question, $U = 0$, represents a geometrical curve such for any given value of the abscissa $t$, the ordinate $z$ has $m$ values: the curve is of the order $m$ at least, but it may be of a higher order $m + \kappa$; when this is so, there is at infinity on the axis $t = 0$, a $\kappa$-tuple point $K$, and thus any line parallel to the axis $t = 0$, intersects the curve in the point $K$ counting as $\kappa$ intersections, and in $m$ other points, which are the points belonging to the $m$ values of the ordinate $z$. Taking $s = 1$, and for $z$, $t$ writing $\dfrac{z}{s}$, $\dfrac{t}{s}$ respectively, we have between the trilinear coordinates $z$, $t$, $s$ an equation $U = 0$ of the order $m + \kappa$, where $s = 0$ is the equation of the line infinity: and the curve has a $\kappa$-tuple point at the point $K$ for which $t = 0$, $s = 0$.

\medskip

\noindent 6.\quad Starting from the equation $U = 0$, for an arbitrary value of $t$ we have $m$ values of $z$ all of them different. These may be developed, each of them as a descending series in $t$, and the development is effected in the usual manner; viz.\ considering for the moment $t$ as very large, then $z$ is in general very large and it has a leading term $At^{\alpha}$, which is determined by writing in the equation $z = At^{\alpha}$, and then finding $\alpha$ in such wise that there are two or more exponents equal to each other and greater (that is, nearer $+\infty$) than any other exponent: we have thus a highest power of $t$, the whole coefficient for which must vanish, viz.\ we obtain an equation of two or more terms giving for $A$ a value or values not $=0$; in the term $At^{\alpha}$ in question, $\alpha$ is in general positive, but it may be $=0$, or be negative. The leading term $At^{\alpha}$ being found in this manner, the law for the exponents of the subsequent terms is frequently at once apparent; but, if necessary, we write $z = At^{\alpha} + A't^{\alpha-\alpha'} + \cdots$, and determine in like manner the exponent $\alpha - \alpha'$ and the coefficient $A'$. Proceeding in this manner until the law of the exponent becomes apparent, and then by the method of indeterminate coefficients, we finally arrive at a series
\[
z = At^{\alpha} + A't^{\alpha-\alpha'} + A''t^{\alpha-\alpha''} + \cdots,
\]
in descending powers of $t$: the exponents are integral or fractional, but the number of positive exponents is always finite. There is either a single series or it may be two or more series; but the number of series is at most $= m$, and when it is less than $m$, then the coefficients in the several series or some of them will contain radicals, and by giving to each radical its different values, the system of series will determine $m$ different values of $z$.

\medskip

\noindent 7.\quad The curve meets the line infinity ($s = 0$) in the point $K$ counting as $\kappa$ intersections, and in $m$ other points, some or all of which may coincide with $K$; the forms of the series depend on the configuration of the $m$ points, or say on the relation of the curve to the line infinity. Thus suppose $\kappa = 0$, and further that the $m$ points are all of them distinct, that is, let the curve be a curve of the order $m$ meeting the line infinity in $m$ distinct points, or, what is the same thing, having $m$ asymptotes, no two of them parallel; we have in this case $m$ series, each of the form
\[
z = At + B + \frac{C}{t} + \frac{D}{t^{2}} + \cdots,
\]
where the several coefficients $A$ have distinct values.

\medskip

\noindent 8.\quad In the foregoing case $A$ is determined by an equation of the order $m$, having $m$ unequal roots; and taking for $A$ any root at pleasure, the remaining coefficients $B$, $C$, \ldots are each of them linearly determined. If the equation has two equal roots, this may correspond to the case of two parallel asymptotes (the curve has here a node at infinity); the coefficients $B$, $C$, \ldots will in this case depend on a quadric radical, and by giving to this radical its two values, we obtain the two series
\[
z = At + B + \frac{C}{t} + \cdots,
\]
corresponding to the two equal roots of the equation. But if instead of a node at infinity we have the line infinity a tangent to the curve, then instead of the two parallel asymptotes, we have an asymptotic parabola; the series assumes a new form, viz.\ it contains terms in $t^{\frac{1}{2}}$, $t^{-\frac{1}{2}}$, \ldots; the coefficients of the integer powers have determinate values, but those of the powers $t^{\frac{1}{2}}$, $t^{-\frac{1}{2}}$, \ldots contain as factor a quadric radical, and giving to this radical its two values, we have two series.

\medskip

\noindent 9.\quad The like considerations apply in the case where the line infinity passes through higher multiple points of the curve, or has with it a contact or contacts other than a single ordinary contact: in every case, reckoning as distinct series those obtained by attaching to each radical its different values, the number of distinct series will be $= m$ precisely.

\medskip

\noindent 10.\quad Reverting now to the differential equation
\[
p_{0}\,[z^{(m-1)} + \cdots + z^{m}] + p_{1}[\cdots + z^{m-1}] + \cdots + p_{m-2}(z' + z^{2}) + p_{m-1}\,z + p_{m} = 0,
\]
it is to be observed that the very same process is applicable to the determination of $z$ in the form of a descending series $z = At^{\alpha} + A't^{\alpha-\alpha'} + \cdots$; moreover, it frequently happens that the determination of the leading term $At^{\alpha}$ is made by means of the algebraical equation
\[
p_{0}\,z^{m} + \cdots + p_{m-2}\,z^{2} + p_{m-1}\,z + p_{m} = 0,
\]
viz.\ this is so whenever the value of $\alpha$ is greater than $1$: in fact, writing $z = At^{\alpha}$ ($\alpha > 1$), then in the sets of terms $\{z^{m} + \cdots + z^{(m-1)}\}, \ldots, (z^{3} + 3zz' + z''), (z^{2} + z')$, the first terms $z^{m}$, $z^{m-1}$, \ldots, $z^{3}$, $z^{2}$ are each of them of a higher degree than any other term in the same brackets, so that attending only to the terms of highest degree the sets may be reduced to $z^{m}$, $z^{m-1}$, \ldots, $z^{3}$, $z^{2}$. But in the determination of the subsequent exponents and coefficients, the omitted terms or some of them would of course come into play: and it is at least not obvious that we can, for the forms of the series which satisfy the differential equation, employ geometrical considerations such as those which were used in regard to the series satisfying the algebraical equation.

\medskip

\noindent 11.\quad Considering as above the coefficients $p_{0}$, $p_{1}$, \ldots, $p_{m}$ as rational and integral functions of $t$, it is easy to determine the degrees of these coefficients in such wise that the exponent $\alpha$ of the leading term $At^{\alpha}$ shall have a given value. Suppose first, this given value is $\alpha = a$, a positive integer greater than $1$: the degrees may be $\theta$, $\theta + a$, $\theta + 2a$, \ldots, $\theta + ma$, ($\theta = 0$ or any positive integer at pleasure); and not only so, but if we write
\[
p_{0} = L_{0}\,t^{\theta} + \cdots,\quad p_{1} = L_{1}\,t^{\theta+a} + \cdots,\quad \ldots,\quad p_{m} = L_{m}\,t^{\theta+ma} + \cdots,
\]
then writing $z = At^{a}$, clearly the highest power in the equation is $t^{\theta + ma}$, and equating the coefficient hereof to zero, we have
\[
L_{0}\,A^{m} + L_{1}\,A^{m-1} + \cdots + L_{m-1}\,A + L_{m} = 0,
\]
an equation of the degree $m$ for the determination of the coefficient of the leading term $At^{a}$ ($a > 1$). Assuming that the roots are all unequal, we have $m$ series each of the form
\[
z = At^{a} + Bt^{a-1} + \cdots + Mt + N + \frac{O}{t} + \cdots.
\]

\medskip

\noindent 12.\quad Suppose $a = 1$, that is, let the leading term be $At$; here the degrees may be $\theta$, $\theta + 1$, $\theta + 2$, \ldots, $\theta + m$, ($\theta = 0$ or any positive integer as before), but we have a different form for the $A$-equation. In fact, here ($z = At$) in the sets of terms $\{z^{m} + \cdots + z^{(m-1)}\}, \ldots, (z^{3} + 3zz' + z''), (z^{2} + z')$, the terms in each brackets are of the same degree, viz.\ the degrees are $m$, \ldots, $3$, $2$; and not only so, but we have $z^{2} + z' = (A^{2} + A)\,t^{2}$, that is, $[A]_{2}\,t^{2}$; $z^{3} + 3zz' + z'' = [A]_{3}\,t^{3}$, \ldots. Hence if
\[
p_{0} = L_{0}\,t^{\theta} + \cdots,\quad p_{1} = L_{1}\,t^{\theta+1} + \cdots,\quad p_{m} = L_{m}\,t^{\theta+m} + \cdots,
\]
the $A$-equation is
\[
L_{0}\,[A]_{m} + L_{1}\,[A]_{m-1} + \cdots + L_{m-1}\,[A]_{1} + L_{m} = 0,
\]
an equation of the degree $m$ as before; assuming that the roots are unequal, we have here $m$ series each of the form
\[
z = At + B + \frac{C}{t} + \frac{D}{t^{2}} + \cdots,
\]
$m$ being, as will presently appear, the case of regular integrals for the $y$-equation.

\medskip

\noindent 13.\quad In either of the last-mentioned cases, the formula for the $A$-equation holds good if any two or more of the coefficients $p_{0}$, $p_{1}$, \ldots, $p_{m}$ are of the degrees aforesaid, the others of them being of inferior degrees; we have only to put $=0$ the $L$'s which belong to the coefficients of inferior degrees. If neither $L_{0}$ nor $L_{m}$ be $=0$, the equation is still an equation of the order $m$, with $m$ effective roots; but if any of the $L$ coefficients at the beginning of the equation vanish we thus introduce roots $=\infty$: and if any of the coefficients $L$ at the end of the equation vanish, we thus introduce roots $=0$: the number of effective roots is $m-$ the number of roots $\infty$ or $0$; the case requires further consideration. If the equation has equal roots, then from what precedes, it is easy to infer that we may have distinct series containing as in the general case only integer powers of $t$, or we may have series beginning with $At^{a}$ or $At$, but containing also fractional powers of $t$.

\medskip

\noindent 14.\quad If the leading term is $z = At^{-a}$, $a$ a value $< 1$, then in the sets of terms $\{z^{m} + \cdots + z^{(m-1)}\}, \ldots, (z^{3} + 3zz' + z''), (z^{2} + z')$ the terms $z^{(m-1)}, \ldots, z''$, $z'$ are each of them of a higher degree than any other term in the same brackets, so that attending only to the terms of the highest degree the sets may be reduced to $z^{(m-1)}, \ldots, z''$, $z'$, or we may consider the equation. In particular, this is the case for $z = At^{-a}$, $a$ being a positive integer; in this case, the degrees of $p_{0}$, $p_{1}$, \ldots, $p_{m}$ may be $\theta + m-1$, $\theta + m$, $\theta + m + 1$, \ldots, $\theta + a$, and this being so, the leading coefficient $A$ will be determined by a linear equation. But the case $\alpha = 0$ is, in fact, that of a non-singular point $x = a$: and I do not here further consider it.

\medskip

\noindent 15.\quad Reverting to the case where the series is
\[
z = At + B + \frac{C}{t} + \frac{D}{t^{2}} + \cdots,
\]
and substituting for $t$ its value $\dfrac{1}{x-a}$, or for greater convenience $\dfrac{1}{x}$, this is
\[
\frac{1}{y}\frac{dy}{dx} = \frac{A}{x} + B + Cx + \cdots,
\]
whence
\[
\log y = A\log x + Bx + \tfrac{1}{2}Cx^{2} + \cdots,
\]
or passing to the value of $y$, but expanding the exponential of $Bx + \tfrac{1}{2}Cx^{2} + \cdots$, this is
\[
y = x^{A}\,(1 + Bx + Cx^{2} + \cdots),
\]
viz.\ we have here a regular integral. By what precedes, it appears that, for the values $p_{0} = L_{0}x^{-\theta} + \cdots$, $p_{1} = L_{1}x^{-\theta-1} + \cdots$, \ldots, $p_{m} = L_{m}x^{-\theta-m} + \cdots$, the exponent $A$ is determined by the equation
\[
L_{0}\,[A]_{m} + L_{1}\,[A]_{m-1} + \cdots + L_{m-1}\,[A]_{1} + L_{m} = 0,
\]
which is the equation for that purpose considered by Fuchs, and is what I would call the indicial equation.

The logarithmic forms which belong to the case of equal roots may be deduced from the case of unequal roots, by supposing two or more of the roots to approach each other continuously and become ultimately equal.

\medskip

\noindent 16.\quad Similarly, if the series be
\[
z = \frac{F}{x^{a}} + \cdots + \frac{K}{x} + A + Bx + \tfrac{1}{2}Cx^{2} + \cdots,
\]
($a$ being a positive integer $=2$ at least), then writing as above $t = \dfrac{1}{x}$, this is
\[
\frac{1}{y}\frac{dy}{dx} = \frac{F}{x^{a}} + \cdots + \frac{K}{x} + A,
\]
whence, if for shortness $w = \dfrac{F}{a-1}\,\dfrac{1}{x^{a-1}} - \cdots - \dfrac{K}{x}$, we have
\[
\log y = w + A\log x + Bx + \tfrac{1}{2}Cx^{2} + \cdots,
\]
or passing to the value of $y$, but expanding as before the exponential of $Bw + \tfrac{1}{2}Cx^{2} + \cdots$, this is
\[
y = e^{w}\,x^{A}\,(1 + Bx + Cx^{2} + \cdots),
\]
a subregular integral. By what precedes, it appears that, if
\[
p_{0} = L_{0}\,x^{-\theta} + \cdots,\quad p_{1} = L_{1}\,x^{-\theta-a} + \cdots,\quad \ldots,\quad p_{m} = L_{m}\,x^{-\theta-ma} + \cdots,
\]
then the equation for the leading coefficient (previously represented by $A$) is
\[
L_{0}\,F^{m} + L_{1}\,F^{m-1} + \cdots + L_{m-1}\,F + L_{m} = 0,
\]
but I do not see any direct or simple way of obtaining the corresponding equation for the exponent $A$ of this subregular integral.

\medskip

\noindent 17.\quad To illustrate the case of a series with fractional powers, and also in order to work out a simple example with some completeness, I consider the differential equation of the second order $\dfrac{d^{2}y}{dx^{2}} + p_{2}\,y = 0$, giving the $z$-equation $z^{2} + z' + p_{2} = 0$; $p_{2}$ is a rational and integral function of $\dfrac{1}{x-a}\,?$, and we have two cases according as the order of the function is even or odd.

\medskip

\noindent 18.\quad Suppose first that the order is even, and for convenience taking it to be $=4$, and assuming the value of $p_{2}$ accordingly, I consider the equation
\[
z^{2} + z' = \frac{\alpha}{x^{4}} + \frac{\beta}{x^{3}} + \frac{\gamma}{x^{2}} + \frac{\delta}{x} + \varepsilon.
\]
Here the form is at once seen to be
\[
z = \frac{A}{x^{2}} + \frac{B}{x} + C + Dx + Ex^{2},
\]
and observing that we have
\[
\frac{dz}{dt} = -\frac{2A}{x^{3}} - \frac{B}{x^{2}} + D + 2Ex,
\]
we find
\begin{align*}
A^{2} &= \alpha, \\
2AB - 2A &= \beta, \\
2AC + B^{2} - B &= \gamma, \\
2AD + 2BC + C &= \delta, \\
2AE + 2BD + C^{2} + D &= \varepsilon, \\
2AF + 2BE + 2CD + 2E &= 0,
\end{align*}
equations which give
\[
A = \pm\sqrt{\alpha},
\]
\[
B = \frac{\beta}{2\sqrt{\alpha}} - \frac{1}{?} = \frac{\beta}{2}\cdot\frac{1}{\sqrt{\alpha}} - 1,\quad ?
\]
\[
C = \frac{\gamma}{2\sqrt{\alpha}} - \frac{1}{2}\cdot\frac{\beta^{2}}{4\alpha} - \cdots,
\]
\[
D = \frac{\delta}{2\sqrt{\alpha}} - \cdots,\qquad E = \frac{?}{?} - \cdots.
\]

\medskip

\noindent We then have
\[
z = \frac{1}{y}\,\frac{dy}{dx} = \frac{A}{x^{2}} + \frac{B}{x} + C + Dx + Ex^{2},
\]
and consequently
\[
\log y = -\frac{A}{x} + B\log x + Cx + \tfrac{1}{2}Dx^{2} + \tfrac{1}{3}Ex^{3} + \cdots,
\]
\[
y = e^{-A/x}\,x^{B}\,e^{Cx + \frac{1}{2}Dx^{2} + \frac{1}{3}Ex^{3} + \cdots},
\]
where the second exponential is to be expanded in the form $1 + Cx + \dfrac{D}{2}x^{2} + \cdots$; there are of course two values, say $y_{1}$ and $y_{2}$, corresponding to the values $+\sqrt{\alpha}$ and $-\sqrt{\alpha}$ of the radical respectively.

\medskip

\noindent 19.\quad Secondly, if the order be odd, and for convenience taking it to be $=3$, and assuming the value of $p_{2}$ accordingly, I consider the equation
\[
z^{2} + z' = \frac{\alpha}{x^{3}} + \frac{\beta}{x^{2}} + \frac{\gamma}{x} + \delta + \varepsilon x.
\]
Here the series is at once seen to be
\[
z = A t^{3/2} + B t + C t^{1/2} + D + E t^{-1/2} + \cdots:
\]
we have
\[
\frac{dz}{dt} = -\tfrac{3}{2}A t^{1/2} - B t^{0} - \tfrac{1}{2}C t^{-1/2} + \tfrac{1}{2}E t^{-3/2} + \cdots,
\]
and we thence derive the equations
\begin{align*}
A^{2} &= -\alpha, \\
2AB - \tfrac{3}{2}A &= 0, \\
2AC + B^{2} - B &= \gamma, \\
2AD + 2BC - \tfrac{1}{2}C &= 0, \\
2AE + 2BD + C^{2} &= \delta, \\
2AF + 2BE + 2CD + \tfrac{1}{2}E &= \varepsilon,
\end{align*}
(where in the next equation we have on the right-hand side $\varepsilon$, and in each of the subsequent equations we have on the right-hand side $0$). We find
\[
A = \pm\sqrt{-\alpha},\quad B = \tfrac{3}{4},\quad C = \tfrac{\gamma}{2A} - \tfrac{?}{?},\quad D = \tfrac{?}{?},\quad E = \tfrac{?}{?},
\]
where observe that the coefficients $A$, $C$, $E$, \ldots contain as a factor $\dfrac{1}{\sqrt{\beta}}$, but the coefficients $B$, $D$, \ldots are rational.

\medskip

\noindent We have
\[
z = \frac{1}{y}\,\frac{du}{dx} = \frac{A}{x^{3/2}} + \frac{B}{x} + \frac{C}{x^{1/2}} + D + Ex^{1/2} + Fx + \cdots,
\]
giving
\[
\log y = -\frac{2A}{x^{1/2}} + B\log x + 2Cx^{1/2} + Dx + \tfrac{2}{3}Ex^{3/2} + \tfrac{1}{2}Fx^{2} + \cdots;
\]
and thence passing to the value of $y$, but expanding the exponential of $2Cx^{1/2} + Dx + \cdots$, we have say the integral
\[
y_{1} = e^{-2A/x^{1/2}}\,x^{B}\,(1 + Cx^{1/2} + Dx + Ex^{3/2} + \cdots),
\]
and then also, changing the sign of $\sqrt{\beta}$, the integral
\[
y_{2} = e^{+2A/x^{1/2}}\,x^{B}\,(1 - Cx^{1/2} + Dx - Ex^{3/2} + \cdots).
\]
Writing for a moment
\[
Q = e^{-2A/x^{1/2}},\qquad M = 1 + Dx + \cdots,\qquad N = C + Ex + \cdots,
\]
the two integrals are $x^{B}\,Q\,(M + N x^{1/2})$, and $x^{B}\,Q^{-1}\,(M - N x^{1/2})$: these may be replaced by their sum and difference, viz.\ these are
\[
y_{1} + y_{2} = M x^{B}\,(Q + Q^{-1}) + N x^{B + 1/2}\,(Q - Q^{-1}),
\]
\[
y_{1} - y_{2} = M x^{B}\,(Q - Q^{-1}) + N x^{B + 1/2}\,(Q + Q^{-1}),
\]
where $Q + Q^{-1}$ is a rational function of $x$, and $Q - Q^{-1}$ is $\sqrt{\alpha}$ multiplied by a rational function of $x$; hence $y_{1} + y_{2}$ is $x^{B}$ multiplied by a rational function of $x$, but $y_{1} - y_{2}$ is $x^{B+\frac{1}{2}}$ multiplied by a rational function of $x$.

\medskip

\begin{flushright}\textit{Cambridge, 28 April, 1886.}\end{flushright}

\vfill
\newpage

\section*{863. Note on the Theory of Linear Differential Equations.}

\begin{flushright}\small
[From \textit{Crelle's Journal der Mathem.}, t.~ci.\ (1887), pp.~209--213.]
\end{flushright}

\medskip

\noindent\textsc{The theorem v.}\ given by Fuchs in the memoir ``Zur Theorie der linearen Differentialgleichungen mit ver\"anderlichen Coefficienten,'' \textit{Crelle's Journal}, t.~\textsc{lxviii}.\ (1868), pp.~354--385 (see p.~374) for the purpose of deciding whether the integrals belonging to a group of roots of the ``determinirenden Fundamentalgleichung'' (or as I call it, the Indicial equation) do or do not involve logarithms, may I think be exhibited in a clearer form.

Starting from the differential equation
\[
P(y) := p_{0}\frac{d^{m}y}{dx^{m}} + p_{1}\frac{d^{m-1}y}{dx^{m-1}} + \cdots + p_{m}y = 0,
\]
of the order $m$, then if $X$ be any function of $x$ not satisfying the differential equation, we can at once form a differential equation of the order $m+1$, satisfied by all the solutions of the differential equation, and having also the solution $y = X$; the required equation is in fact
\[
d_{x}P(y)\cdot P(X) - P(y)\cdot d_{x}P(X) = 0.
\]
This I call the augmented equation.

I recall that the equation $P(y) = 0$, considered by Fuchs, is an equation having for each singular point $x = a$, $m$ regular integrals, viz.\ the coefficients $p_{0}$, $p_{1}$, \ldots, $p_{m}$, have the forms $q_{0}(x-a)^{m}$, $q_{1}(x-a)^{m-1}$, \ldots, $q_{m}$, where $q_{0}$, $q_{1}$, \ldots, $q_{m}$ are rational and integral functions of $x-a$, $q_{0}$ not vanishing for $x = a$, and the other functions $q_{1}$, $q_{2}$, \ldots, $q_{m}$ not in general vanishing for $x = a$. Writing $y = (x-a)^{\theta}$, we obtain
\[
P(x-a)^{\theta} = q_{0}\,I(\theta)\,(x-a)^{\theta} + \text{higher powers of }(x-a),
\]
where $I(\theta)$, the coefficient of the lowest power of $(x-a)$, is a function of $\theta$ of the order $m$, which I call the indicial coefficient; and equating it to zero, we have $I(\theta) = 0$, the determinirende Fundamentalgleichung, or Indicial equation, being an equation of the order $m$. If the roots of this equation are such that no two of them are equal or differ only by an integer number, then we have $m$ particular integrals each of them of the form
\[
y = (x-a)^{r} + \text{higher powers of }(x-a),
\]
where $r$ is any root of the indicial equation: but if we have in the indicial equation a group of $\lambda$ roots $r_{1}, r_{2}, \ldots, r_{\lambda}$, such that the difference of each two of them is either zero or an integer, then the integrals which correspond to these roots involve or may involve logarithms; in particular, if any two of the roots are equal, the integrals for the group will involve logarithms.

Consider now the differential equation $P(y) = 0$ in reference to the singular point $x = a$ as above, and writing $X = (x-a)^{e}\,f$ where $e$ is in the first instance arbitrary, and $f$ is a rational and integral function of $x-a$ not vanishing for $x = a$, we form the augmented equation which, observing that we have in general $P(X) = (x-a)^{e}\,Q$, $Q$ a rational and integral function of $x-a$ not vanishing for $x = a$, and dividing the whole equation by $(x-a)^{e-1}$, may be written
\[
d_{x}P(y)\cdot (x-a)Q - P(y)\,\bigl[e\,Q + (x-a)\,d_{x}Q\bigr] = 0,
\]
an equation of the same form as the original equation (but of the order $m+1$ instead of $m$), and having an indicial equation
\[
(\theta - e)\,I(\theta) = 0.
\]
In fact, writing as before $y = (x-a)^{\theta}$, we have in $d_{x}P(y)\cdot (x-a)\,Q$ the term of lowest order $\theta\,I(\theta)\,Q_{0}\,(x-a)^{\theta-e}$, and in $P(y)\cdot e\,Q$ the term of lowest order $e\,I(\theta)\,Q_{0}\,(x-a)^{\theta-e}$, whereas in $P(y)\,(x-a)\,d_{x}Q$ the term of lowest order is $(x-a)^{\theta-e+1}$; the indicial equation is thus as just found.

If however $e$ be equal to a root of the indicial equation $I(\theta) = 0$, then instead of $P(X) = (x-a)^{e}\,Q$, we have $P(X) = (x-a)^{\mu}\,Q$, where the index $\mu$ is $= e + a$ positive integer, and where the value of the difference $\mu - e$ may depend upon the determination of the function $f$ in the expression $(x-a)^{e}\,f$. The indicial equation for the augmented equation is in this case $(\theta - \mu)\,I(\theta) = 0$.

If the indicial equation $I(\theta) = 0$ of the given differential equation has a group of roots $r_{1}, r_{2}, \ldots, r_{\lambda}$, the difference of any two of these roots being zero or an integer, then taking $e = $ any one of these roots, the augmented equation will have a group of roots $(\mu, r_{1}, r_{2}, \ldots, r_{\lambda})$.

If any two of the roots $r_{1}, r_{2}, \ldots, r_{\lambda}$ are equal, the group of integrals $u_{1}, u_{2}, \ldots, u_{\lambda}$ will involve logarithms: the question only arises when these roots are unequal, and taking them to be so, the theorem~v.\ is in effect as follows: ``If by taking $e = $ some one of the roots $r_{1}, r_{2}, \ldots, r_{\lambda}$, and by a proper determination of the function $f$ we can make $\mu$ to be $= $ one of the same roots $r_{1}, r_{2}, \ldots, r_{\lambda}$, then the group of integrals $u_{1}, u_{2}, \ldots, u_{\lambda}$ will involve logarithms; but if $\mu$ cannot be made $= $ one of the roots $r_{1}, r_{2}, \ldots, r_{\lambda}$, then the group of integrals will be free from logarithms.''

As an example, I consider the equation
\[
\frac{d}{dx}\!\left[(1-x^{2})\,\frac{d^{2}y}{dx^{2}}\right] - 2x\,\frac{dy}{dx} + n(n+1)\,y = 0.
\]
This is Legendre's equation $(1-x^{2})\,\dfrac{d^{2}y}{dx^{2}} - 2x\,\dfrac{dy}{dx} + n(n+1)\,y = 0$, with $\dfrac{1}{x}$ substituted for $x$, so that, instead of a singular point $x = \infty$, there may be a singular point $x = 0$. Attending to the singular point $x = 0$, we have
\[
P(x)^{\theta} = \bigl[(\theta - n)(\theta - n - 1) - 2\theta\bigr]\,x^{\theta} + \text{higher powers},
\]
so that the indicial equation $I(\theta) = 0$ is $\theta^{2} - \theta - n^{2} - n = 0$, that is, $(\theta + n)(\theta - n - 1) = 0$, or we have the roots $-n$, $n + 1$, which differ by an integer, and thus form a group, if $n$ be $= $ an integer, or be $= $ an integer $+ \tfrac{1}{2}$; to fix the ideas, say that the roots are $-p$, $p + 1$ or else $-p - \tfrac{1}{2}$, $p + \tfrac{1}{2}$, where $p$ is a positive integer.

Writing for greater convenience $x^{e}f = x^{e} + F$, where $F$ is a sum of powers of $x$ higher than $e$, we find without difficulty
\[
P\bigl(x^{e}f\bigr) = x^{e}\,\bigl\{(e-n)(e-n-1) - 2e\bigr\}\,(x^{0} - x^{-2} + x^{4})\,F^{?} + (n? \cdot n)\,x^{-e}\,F\,?
\]
which, so long as $e$ remains arbitrary, is of the form $x^{e}\,Q$, $Q = (e + n)(e - n - 1) + $ powers of $x$; if however $e$ be a root of the indicial equation, $=$ for instance, if $e = -n$, then the expression in brackets $\{\ \}$ contains at any rate the factor $x$, so that the form is
\[
P\bigl(x^{-n}f\bigr) = x^{\mu}\,Q,
\]
where $\mu$ is $= -n + 1$ at least; we can however, by a proper determination of the function $f$, make $\mu$ acquire a larger value.

For instance, suppose $-n$, $n + 1 = -2$, $3$, $\therefore$ $n = 2$, and assume
\[
f = x^{e} + x^{-2} + Bx^{-1} + Cx^{0} + Dx^{1} + Ex^{2} + Fx^{3} + Gx^{4} + \cdots
\]
To calculate $P\bigl(x^{e}f\bigr)$, we have

\medskip

\textit{[Calculation table follows: the polynomial $P\bigl(x^{e}f\bigr)$ is set out in rows giving the coefficients of $x^{e-2}$, $x^{e-1}$, $x^{e}$, $x^{e+1}$, \ldots, $x^{e+5}$, etc., as a function of the indeterminates $A$, $B$, $C$, $D$, $E$, $F$, $G$, and the parameter $\mu = e + n$.]}

\medskip

\noindent Hence if $B$ not $=0$, we have $\mu = -1$; if $B = 0$, $C$ not $=0$, we have $\mu = 0$; if $B = 0$, $C = 0$, $D$ not $=0$, we have $\mu = 1$; if $B = 0$, $C = 0$, $D = 0$, $E$ not $=0$, we have $\mu = 2$; if $B = 0$, $C = 0$, $D = 0$, $E = 0$, $F$ arbitrary, we can by giving proper values to the higher coefficients $B$, $C$, $D$, $E$, $\ldots$ from this on make $\mu$ have any integer value. The values of $\mu = -1$, $0$, $1$, $2$ being all of them less than the indicial roots $-p$, $p + 1$, we thus see by the theorem that, $E$ not $=0$, then by a proper determination of the function $f$ in the equation $X = x^{e}\,f$, we can bring $\mu$ up to the value $p$, and this being so, the indicial equation of the augmented equation is $\bigl(\theta - p\bigr)(\theta + p)(\theta - p - 1) = 0$; that is, the group of integrals $u_{1}, u_{2}, \ldots, u_{p}$ which correspond to the roots $-p$, $p + 1$ of the original equation will involve logarithms. The case $B = 0$, etc.\ is one which appears to me to require further consideration; but I do not here further consider it.

\vfill
\end{document}
